\section{Approche réflexive}

\subsection{Les connaissances/compétences acquises grâce à la rédaction de ce cas clinique ;} 

En cours de cette hospitalisation, j’ai pu prendre le temps de comprendre en profondeur le dossier de monsieur. Rechercher les informations pertinentes parmi tous les comptes-rendus disponibles dans son dossier. Synthétiser un cas aussi complexe est déjà un exercice en lui-même.\\
De plus, vient s’ajouter l’aspect pharmacologique. Il faut comprendre la complexité des traitements médicamenteux et leur impact concret sur le patient et sa clinique. Remettre en question la légitimité de chaque molécule. Chez Monsieur, je suis convaincue que cette tâche n’a pas été exploitée dans son entièreté. Est-ce que la vitamine D et le traitement pour l’hypercholestérolémie avaient réellement leur place dans le pilulier de Monsieur vu le contexte palliatif ? Ou encore le Zolpidem, une Z-drogue que l’on doit remettre en question systématiquement vu l’usage abusif à long terme de cette classe de médicament. Son indication est à remettre en question  mais le sevrage n’est peut-être pas à envisager en aigu lors de son hospitalisation et dans le contexte des soins palliatifs.

\subsection{Les influences probables de ce travail sur sa pratique médicale future ;}

L’importance de clarifier le projet de soin et de façon précoce dans l’hospitalisation afin de prévenir une mauvaise gestion d’événement aigu. Au-delà de la connaissance médicale, notre rôle est d’aider le patient et de le soulager. Pour pratiquer de la bonne médecine, il faut être en accord avec le patient et considérer le patient comme l’expert de son propre corps.\\ 

La communication peut participer activement à la prise en soin du patient. Les enjeux sont plus grands que ce que la plupart des médecins veulent bien admettre.

\subsection{Les difficultés (dépendantes ou indépendantes de sa volonté) ayant entravé son acquisition de nouvelles connaissances/compétences au cours de ce travail ;}

Les obstacles rencontrés autour de ce cas ont été l’accessibilité et la flexibilités des acteurs extérieurs à l’hôpital afin de coordonner la sortie du patient. Étant dans la campagne de la province du Luxembourg, la disponibilité des soins est réduite. Le médecin généraliste en charge a joué un rôle majeur dans l’organisation des soins à domicile. Dans un premier temps, il a été compliqué de joindre celui-ci. Heureusement, il s’est montré très flexible par la suite mais cela nous rappelle à quel point les soins sont parfois inaccessible en campagne.\\

Étant proche de la frontière luxembourgeoise,  la comparaison des remboursement disponible entre les deux pays est aisée. Une fois de plus, les possibilités de soins sont différentes. Parfois, pour un même profil de patient, la prescription et la qualité de soins seront très différentes selon les assurances respectives des patients.\\

D’un point de vue médical, devoir contrôler tous les facteurs et trouver l’équilibre bénéfices/ risques des traitements qui étaient instaurés. De plus, le manque de temps et l’urgence qui s’ajoutait au tableau clinique. 

\subsection{Les solutions que l’étudiante peut mettre en place pour remédier à ces difficultés ;}

Persévérer et tenter d’anticiper un maximum les démarches administratives afin d’avoir une chance de pouvoir les mettre en place. Cela peut être assez simple à mettre en place en commençant par s’entourer des bonnes personnes, contacter sans tarder l’équipe de soins continus et sociale le plus vite possible. Les qualités indispensables à un bon (stagiaire) médecin sont d’avoir l’esprit de groupe et de profiter des ressources de l’hôpital. La complexité de ces cas requiert l’avis de plusieurs spécialistes, une médecine solitaire n’est pas optimale et est à proscrire. 
